\section{The interaction picture in Hastings' approximate ground-state projector}
\label{sec:hastings_interaction_picture}

% Hastings, arXiv:0705.2024v4, law16.tex, lines 843--901 (step five of the
% proof of Lemma 1). Finite complex matrices; no lattice content.

Step five of the proof of \cite[Lemma~1]{Hastings2007AreaLaw} factors
$e^{i(A+B)t}=U(t)e^{iAt}$ with $A=M_L+M_R$ and $B=M_B$.

\begin{definition}[Interaction-picture propagator]
    \label{def:hastings_interaction_propagator}
    \lean{TNLean.MPS.AreaLaw.Hastings.interactionGenerator,
        TNLean.MPS.AreaLaw.Hastings.interactionPropagator}
    \leanok
    For square matrices $A,B$ and $t\in\R$, set
    \begin{align}
        B(t) & =e^{iAt}Be^{-iAt}, & U(t) & =e^{i(A+B)t}e^{-iAt}.
        \label{eq:hastings_interaction_propagator}
    \end{align}
\end{definition}

\begin{lemma}[Interaction-picture factorization]
    \label{thm:hastings_interaction_picture}
    \lean{TNLean.MPS.AreaLaw.Hastings.interactionPropagator_hasDerivAt,
        TNLean.MPS.AreaLaw.Hastings.eq_interactionPropagator_of_hasDerivAt,
        TNLean.MPS.AreaLaw.Hastings.interactionPropagator_mem_unitaryGroup,
        TNLean.MPS.AreaLaw.Hastings.hermitianUnitaryPath_add_eq_interactionPropagator_mul}
    \leanok
    \uses{def:hastings_interaction_propagator}
    The path $U$ is the unique differentiable solution of
    \begin{align}
        U(0)=1,\qquad U'(t)=U(t)\,iB(t),
        \label{eq:hastings_interaction_equation}
    \end{align}
    satisfies
    \begin{align}
        e^{i(A+B)t}=U(t)e^{iAt},
        \label{eq:hastings_interaction_factorization}
    \end{align}
    and is unitary when $A=A^\dagger$ and $B=B^\dagger$
    \cite[lines 853--864]{Hastings2007AreaLaw}.
\end{lemma}
\begin{proof}\leanok
    Differentiating~(\ref{eq:hastings_interaction_propagator}) gives
    (\ref{eq:hastings_interaction_equation}) and
    (\ref{eq:hastings_interaction_factorization}). For uniqueness,
    $V(t)=e^{iAt}e^{-i(A+B)t}$ satisfies $V'(t)=-iB(t)V(t)$ and $V(t)U(t)=1$,
    so any solution $W$ of~(\ref{eq:hastings_interaction_equation}) has
    $(WV)'=0$, hence $W(t)V(t)=1$ and $W=U$.
\end{proof}

\begin{lemma}[Cut-off commutation estimate]
    \label{thm:hastings_cutoff_commutation}
    \lean{TNLean.MPS.AreaLaw.Hastings.norm_hermitianUnitaryPath_mul_sub_le}
    \leanok
    For $K=K^\dagger$, any $X$ of the same size, and $t\in\R$,
    \begin{align}
        \|e^{iKt}X-X\|\leq|t|\,\|KX\|
        \label{eq:hastings_cutoff_commutation}
    \end{align}
    (applied with $K=M_L+M_R$ and $X=O_LO_R$ in
    \cite[lines 878--880]{Hastings2007AreaLaw}).
\end{lemma}
\begin{proof}\leanok
    Since $\frac d{ds}(e^{iKs}X)=e^{iKs}\,iKX$ has norm $\|KX\|$, the mean
    value inequality on $[0,t]$ gives~(\ref{eq:hastings_cutoff_commutation}).
\end{proof}

\begin{definition}[Gaussian average of the propagator]
    \label{def:hastings_interaction_average}
    \lean{TNLean.MPS.AreaLaw.Hastings.interactionAverage,
        TNLean.MPS.AreaLaw.Hastings.interactionAverage_eq_integral}
    \leanok
    \uses{def:hastings_interaction_propagator}
    For $h\ge0$ and the centered Gaussian measure $\gamma_h$ of variance $h$,
    set
    \begin{align}
        P_B=\int U(t)\,d\gamma_h(t)
        \label{eq:hastings_interaction_average}
    \end{align}
    (for $h=q/\Delta E^2>0$, the operator $P_B$ of
    \cite[lines 891--897]{Hastings2007AreaLaw}).
\end{definition}

\begin{lemma}[The Gaussian average is a contraction]
    \label{thm:hastings_interaction_average_norm}
    \lean{TNLean.MPS.AreaLaw.Hastings.norm_interactionAverage_le}
    \leanok
    \uses{def:hastings_interaction_average}
    If $A=A^\dagger$ and $B=B^\dagger$, then $\|P_B\|\leq1$ for every $h\ge0$.
\end{lemma}
\begin{proof}\leanok
    \uses{thm:hastings_interaction_picture}
    By Lemma~\ref{thm:hastings_interaction_picture}, $\|U(t)\|=1$, and
    $\gamma_h(\R)=1$.
\end{proof}

\begin{lemma}[Moving the cut-off past the boundary evolution]
    \label{thm:hastings_interaction_average_last}
    \lean{TNLean.MPS.AreaLaw.Hastings.norm_interactionAverage_mul_sub_le,
        TNLean.MPS.AreaLaw.Hastings.norm_interactionAverage_mul_mul_sub_le}
    \leanok
    \uses{def:hastings_interaction_average}
    For $A=A^\dagger$, $B=B^\dagger$, matrices $X,P_0$, and $h\geq0$,
    \begin{align}
        \Bigl\|P_BX-P_0\Bigr\|
        \leq\Bigl\|\int e^{i(A+B)t}\,d\gamma_h(t)\,X-P_0\Bigr\|
        +\sqrt h\,\|AX\|.
        \label{eq:hastings_interaction_average_last}
    \end{align}
    For $A=M_L+M_R$, $B=M_B$, $X=O_LO_R$, and $h=q/\Delta E^2$, this gives
    \cite[eq.~(last), lines 881--890]{Hastings2007AreaLaw}.
\end{lemma}
\begin{proof}\leanok
    \uses{thm:hastings_interaction_picture,thm:hastings_cutoff_commutation,
        thm:hastings_gaussian_absolute_moment}
    By~(\ref{eq:hastings_interaction_factorization}) and
    (\ref{eq:hastings_cutoff_commutation}),
    $\|e^{i(A+B)t}X-U(t)X\|=\|U(t)(e^{iAt}X-X)\|\leq|t|\,\|AX\|$. Integrating
    against $\gamma_h$ via Lemma~\ref{thm:hastings_gaussian_absolute_moment}
    gives $\|\int e^{i(A+B)t}\,d\gamma_h\,X-P_BX\|\leq\sqrt h\,\|AX\|$, and the
    triangle inequality yields~(\ref{eq:hastings_interaction_average_last}).
\end{proof}
